\documentclass[11pt]{article}
\usepackage{amsmath,amssymb,amsthm,geometry}\geometry{margin=1in}
\newtheorem{proposition}{Proposition}
\title{Atomic polynomial closure requires entry and external-response structure}
\date{17 November 2026}\begin{document}\maketitle
\paragraph{Scope and target.} This is a discriminating probe for the proposed independent progressive polynomial action on globally boundary-labelled atomic protocols. It tests the direct extension of the reviewed $P_C$ on categories (\texttt{Progressive\_Graded\_Spans\_and\_Causal\_Lift.tex}), not every possible enriched polynomial action. The existing synchronized-history theorem (\texttt{Synchronized\_History\_Activation\_and\_Partial\_Covers.tex}) assumes a common atomic input interface; it does not construct that interface. The present counterexamples identify additional fields needed to make the polynomial itself respect actual activation. They use only finite deterministic data, so probability or coefficient refinements cannot repair them alone.

\section{Candidate input class and preservation contract}
Suppose a base graph records primitive actor states and addresses, physical boundary labels, selected paid edges and actual passive edges. Passive edges in the pinned activation doctrine preserve the actor's live occurrence observation. New actors may be installed only by the supplied initializer interface, at the actual postworld, and cannot take a paid edge in their birth scope. A prospective $P_C^{\rm at}$ should agree with the independent content $P_C$ after forgetting these decorations, have full unused grade $C$, and preserve these operational restrictions. The strong forgetting requirement means retaining all the objects and arrows of the ordinary $P_C$, not only mapping into it. The examples below refute this strong requirement on entry-sensitive activation data.

\section{An ordinary progressive injection can install a non-entry state}
Take terminal world and one unary component $M$. Its actual initial raw carrier is the singleton $z_0$, its initializer law is forced, its selected state is $s_0$, and its initial context coefficient is the unit. A distinct state $s_1$ is reached only after a paid edge $s_0\to s_1$. No silent entry path reaches $s_1$. Both states may have the same output label and same declared full terminal grade; state identity and paid history remain part of the interface.

Let $B$ be the discrete category on $s_0,s_1$. For every collection $C$, the ordinary progressive category $P_CB$ contains
\[
                         u()\longrightarrow u(s_1).       \tag{1}
\]
Its position injection is empty, so it requires no $B$ arrow at all. If (1) is interpreted as an atomic port birth installing $M$ in state $s_1$, it violates the supplied initial selected-state field. It cannot be implemented by first installing $s_0$ and then executing the paid edge: that is forbidden newborn activity in the birth scope, and would add a paid event and context factor.

\begin{proposition}[Failure of entry-blind atomic polynomial preservation]
There is no atomic refinement which both retains every ordinary $P_C$ arrow on this object set and interprets every new-position injection as installation of the supplied static component with its actual initial raw/selected-state interface, while preserving paid events and forbidding newborn paid activity.
\end{proposition}
\begin{proof}
The retained arrow (1) would have to install $s_1$ through the only initializer witness $z_0$, whose selected state is $s_0\ne s_1$. The only supplied alternative path includes a paid edge and is excluded by the preservation contract.\end{proof}

An event source can be supplied externally to make the ambient scope nonempty; the state mismatch persists. Initializing a different component with seed $s_1$, adding a silent reset, or quotienting $s_0,s_1$ changes the given interface. Restricting $P_C$ to injections whose new objects are genuine entry-selected states is a viable direction, but is a changed arrow class and needs its own composition and multiplication proof. In particular, sequential entry then paid evolution is not equivalent to a direct atomic birth at its final state.

\section{External response is not port-invariant passivity}
Take an existing downstream activation bundle with no currently live component. An upstream old actor executes a paid edge exposing its first $B$ port. Choose $M$ so that its initializer-selected entry immediately reveals one live output. Actual owner-aware activation installs that fresh $M$ entry in the downstream bundle at the common postworld, changing its visible output list from empty to nonempty. The downstream bundle contributes no paid primitive edge in this scope, but its actor/port configuration changes.

If a proposed polynomial graph class offers only two classifications for every old composite component, namely ``contains a local paid edge'' or ``actual port-invariant passive response'', this downstream input-response step belongs to neither. It has no local paid primitive actor, so the first classification is false. It changes the visible output list from empty to nonempty, so treating it as the pinned port-invariant passive response is false. The upstream paid actor is the actual owner/sponsor of the entry; it must not be copied into the downstream component merely to make a local activity flag true.

\begin{proposition}[Two-mode local classification is not closed under activation input]
A class with only those two local edge modes cannot represent this actual upstream-driven installation while preserving primitive actor identity, passivity and once-owned costs.
\end{proposition}
\begin{proof}
The two available modes contradict, respectively, the absence of a downstream paid actor and the change in installed downstream configuration. Installing a fictitious active actor or charging entry to a new local paid event violates the stated preservation requirements.\end{proof}

This is a statement about a whole response bundle, not a claim that an old passive primitive actor itself gives birth. Every old passive primitive remains passive; the changed bundle gains a new actor in response to the upstream raw observation. A typed external-response edge with an explicit sponsor/owner boundary is the missing third mode. It may be combined with old local paid activity, but must remain distinguishable from both that activity and primitive passivity.


\section{Why a common boundary word alone does not close entry restriction}
Retaining a common world or atomic-word label, including at nullary tuples, is necessary but not sufficient if a newborn position in every arrow must be a designated entry at that arrow's final world. Let the component have entry $s_0$ and a paid transition $s_0\to s_1$ as above. An external sponsored atomic step first births $s_0$. A second atomic step lets this now-old component execute its paid transition to $s_1$. Both steps satisfy the local entry and no-newborn-paid rules. Their composite has a position absent at its source and present as $s_1$ at its target. Thus it fails the same final-entry-only predicate, even when the composite arrow retains the entire two-letter boundary word and the nullary source world.

This is not an objection to a free-path construction: the two-letter path is valid and remembers that the actor was born in its first letter and acted only in the second. It is an obstruction to restricting \emph{all category arrows} by the atomic final-entry predicate. A closed category must instead distinguish atomic generators from paths and retain birth-time/suffix information or full intermediate configurations. Ordinary word labels without per-position birth history do not establish that information. With paid identity-world updates, world equality alone is even weaker and cannot recover the number of scopes.

An enriched fibrewise construction can therefore use a common atomic word and, for each position, its first-birth index plus its actual subsequent component path. Newborn composite entries must expand to their component entries at that same index; an inner component born earlier may be active later but never in its birth letter. A synchronized graded multiplication would need to flatten these records at every letter, not just at final endpoints. It must also transport actual external sponsor/request annotations and preserve the entry class. These are explicit candidate data for a positive construction; the ordinary $P_C$ and its Cat-over-world variant without such records do not supply them automatically.

\section{What a corrected polynomial action must retain}
The proposed repair must therefore have an independently typed entry interface: either distinguished deterministic entries, or formal pending requests with every admitted raw entry witness and its selected-state map. The actual initializer law families and context coefficients remain separate operational fields. A newborn composite entry must include the entries at its already revealed child ports, including all raw alternatives and zero branches. Entry closure under $P_C$ and under its multiplication is a theorem obligation, not a consequence of ordinary tuple concatenation.

Its atomic boundary must also expose the external sponsor and primitive owner/request disclosure graph. A component response may install actors without contributing a local paid edge; its cost remains in the producing actor's whole expanded factor. Same-event generators must preserve the globally old primitive actor set, chosen old paid edges and silent newborn paths. The synchronized associator proves regrouping only after these annotations are present. It does not infer them from final output lists.

A useful distinction limits the obstruction. If polynomial grouping nodes are genuinely \emph{content-only}, they add no primitive actors. Then flattening a list of lists preserves all leaf actor objects, even when a leaf actor has no revealed output. Loss of an empty intermediate grouping node is not by itself loss of an actor. Conversely, representing an unfinished activation component solely by its exposed output list erases a real actor whenever that list is empty. Such an encoding cannot be used to establish atomic preservation. The base must say which case it represents.

Ordinary progressive multiplication can additionally forget latent phase/shape/grouping data before completion. The graded-span apex retains that data and closes content composition, but no span arrow thereby acquires an entry law or an executable external response. Entire unused grades and actual content arrows remain independent of the owner-context factors, including in these unit-cost counterexamples.

\paragraph{Exact next obligation.}
An entry-generated, sponsor-labelled atomic graph class may evade both counterexamples by restricting ordinary injections and adding typed response generators. To claim closure it must define its pure maps, its composite-entry construction, $\eta$ and $\mu$ on every annotated generator, and show the same admissible joint request family and actual owner maps along the mixed faces. Pullbacks and noninjective pure input maps must retain the sponsor/request indices rather than merge them without a compatibility square. No such completed enriched-polynomial theorem is asserted here. The two propositions rule out the direct entry-blind/two-mode candidate and identify precisely why the existing Cat content algebra is not already that theorem.
\end{document}
